\begin{beginnerstatement}[Kurz erklärt: Continuous Integration]

\emph{Continuous Integration} (CI) bedeutet, dass Prüfungen nach
Quellcodeänderungen automatisch auf einem separaten System ausgeführt werden.
Im untersuchten Projekt übernimmt GitHub Actions diesen Dienst. Ein
\emph{Workflow} beschreibt einen solchen automatischen Ablauf. Darin ist ein
\emph{Job} eine einzelne Arbeitsgruppe von Schritten, die auf einem
\emph{Runner}, also dem bereitgestellten Rechner, läuft. Eine
\emph{Pipeline} ist die gesamte Kette dieser automatischen Prüfungen. Ein
Commit bezeichnet einen festen Codezustand; \emph{HEAD} zeigt auf den Commit,
der für die betrachtete Arbeitskopie oder den Zweig gerade maßgeblich ist.
Eine grüne CI hat alle verlangten Jobs bestanden, eine rote CI enthält
mindestens einen Fehlschlag. Die im
Haupttext ausgewerteten Läufe für \texttt{a4487ac} und
\texttt{461fc751} sind commitgebundene Momentaufnahmen. Beim neueren Stand war
der Master-Quality-Job grün, mehrere Profil- und Windows-Pfade aber nicht.
Ein grüner Job macht deshalb nicht automatisch die gesamte CI grün.

\end{beginnerstatement}
